\section{Background}
Whilst \ce{CO_2} emissions are often portrayed as the leading contributor to aviation-induced climate change, they are only responsible for a mere third of the climate impact. The remaining two-thirds are attributed to more reactive non-\ce{CO_2} emissions, \ce{NO_x}, \ce{H_2O} and PM \cite{Lee2021}. These emission species interact with ambient air through chemical and microphysical processing, giving rise to the production and depletion of radiatively active substances that perturb the net energy balance of the atmosphere (e.g. \ce{NO_x}-induced ozone production, contrail generation through \ce{H_2O} and PM emissions etc.). The degree to which aircraft emissions induce a climatic response varies depending on the state of the background atmosphere (i.e. its chemical composition and meteorology) and the time of day and year in which emissions are released. This means that aviation climate impact is spatio-temporally sensitive, i.e. the same emissions released at different times and/or locations can lead to very different climate effects. 

The dispersion of aircraft emissions occurs over great distance and time scales, with emissions entrained in the aircraft exhaust plume which spreads hundreds of kilometres over its lifetime of up to 12 hours \cite{Kraabol2000a, EPA1992}. The elevated concentrations of emitted chemical species present within the plume result in additional nonlinear chemical (gas-phase and heterogeneous) and microphysical processes which are not accounted for in global chemistry models, due to the inherent assumption of instantaneous dispersion (ID) of emissions. The ID assumption models emissions as homogeneously mixed into the volume of the computational grid cell to which they are released, thus neglecting any subgrid scale nonlinear processing that may occur throughout the plume lifetime \cite{Paoli2011}. Plume-scale nonlinear effects are also further augmented in high-density airspace regions, where plumes intersect and the emissions contained within them accumulate, leading to chemical saturation. Two key saturation effects have been noted to be of particular significance with respect to aviation climate impact: (1) the saturation of \ce{NO_x} emissions, which leads to decreasing ozone production with increasing \ce{NO_x} concentrations \cite{Jaegle1999}, and (2) the dehydration of water vapour leading to diminished contrail climate impact \cite{Schumann2015}. Therefore, the atmospheric response to non-\ce{CO_2} aircraft emissions is not only sensitive to natural variations in the atmospheric state, but it is also affected by heightened emissions concentrations due to the presence of lingering plumes from previous aircraft.

Optimising flight routes with respect to minimum climate impact (instead of minimum fuel burn) is a concept that has become prevalent in the literature in recent years. This involves re-routing aircraft to avoid particularly climate-sensitive regions of the atmosphere, based on provision of en-route chemical and meteorological information. Nikla{\ss} et al.\ \cite{Niklass2019a} have shown through simulation efforts that it is possible to achieve a 12\% reduction in climate impact at virtually no additional cost to fuel burn, evidencing the practicability of climate-optimised routing. Other studies approach this issue purely from the perspective of superimposing aircraft plumes through formation flight. The aerodynamic benefits of formation flight due to wake energy retrieval are already well established \cite{Bangash2004, Kent2020}, however it has become increasingly evident that formation flight can also be used to exploit climate-beneficial saturation effects due to the superposition of plumes from aircraft involved. In Dahlmann et al.\ \cite{Dahlmann2020} a preliminary study exploring the potential for climate impact reduction due to saturation effects in formation flight and found that \ce{NO_x} saturation led to a reduction in ozone production efficiency of \textasciitilde5\% and a mutual inhibition of contrail growth, quelling their radiative properties by 20 to 60\%. 

This review chapter outlines the current state of literature on aviation's atmospheric effects and the influence of nonlinear plume-scale processing on the global net climate impact, to bring to light the notion that aviation-induced perturbations to atmospheric chemistry can simply be viewed as another constraint in the climate-optimised routing problem. Section \ref{emissions} covers the emissions generation process and the methods used to model aircraft performance and emissions. Section \ref{dispersion_2.3} explores the dispersion characteristics of aircraft emissions on the subgrid scale, alluding to the various modelling methods developed to investigate the dynamical characteristics of aircraft plumes. In section \ref{airtraf_2.4}, air traffic management principles and their effect on aviation emissions distribution is considered, on both the local and global scale, followed by section \ref{climate} which explores aircraft climate impact, based on the emissions released and the chemical and physical processes that occur on the grid and subgrid scale of global models. Section \ref{2.6} looks into the effect of subgrid scale processes that occur due to emissions accumulation in high-density airspace, and finally section \ref{Mitigation} explores potential mitigation efforts to reduce aviation-induced climate change by modifying aircraft operations, such as climate-optimal routing and formation flight.